\documentclass[12pt,a4paper]{article}
\usepackage[utf8]{inputenc}
\usepackage{amsmath, amssymb, amsthm}
\usepackage{mathtools}
\usepackage{geometry}
\geometry{margin=1in}

\newtheorem{theorem}{Theorem}[section]
\newtheorem{lemma}[theorem]{Lemma}
\newtheorem{definition}[theorem]{Definition}
\newtheorem{remark}[theorem]{Remark}

\title{On the Iota Function and the Greedy Condition: \\ A Reformulation, Not a Proof}
\author{}
\date{}

\begin{document}

\maketitle

\begin{abstract}
The invertible geometric encoding \(\iota\) maps a complex number \(s\) to a spiral in \(\mathbb{R}^3\) whose horizontal projections are the terms of the Dirichlet eta series. We show that the greedy condition (GC) for a zero of \(\eta(s)\) can be expressed as a geometric property of this spiral. However, proving that every zero satisfies GC remains equivalent to the Riemann Hypothesis. The iota function does not provide an unconditional proof; it merely offers a new visualisation. We present the equivalence and explain why no unconditional derivation of GC from \(\eta(s)=0\) exists without assuming RH.
\end{abstract}

\section{Introduction}

The iota function \(\iota(s)\) was introduced in the original summary as an injective map from \(\mathbb{C}\) to a sequence of points in \(\mathbb{R}^3\):
\[
\iota(s) = (P_1,P_2,\dots),\quad 
P_n = \bigl(n^{-\sigma}\cos\Phi_n,\; n^{-\sigma}\sin\Phi_n,\; n\bigr),\quad 
\Phi_n = \pi(n-1)-t\ln n\ (\text{unwrapped}).
\]
The horizontal projection of \(P_n\) is the vector \(\mathbf{v}_n = (-1)^{n-1}n^{-s}\). The partial sums \(\mathbf{S}_N = \sum_{n=1}^N \mathbf{v}_n\) satisfy \(\mathbf{S}_N \to 0\) if and only if \(\eta(s)=0\).

The greedy condition (GC) is
\[
\langle \mathbf{S}_{N-1},\; \mathbf{v}_N \rangle \le 0 \qquad \forall N\ge 1,
\]
where \(\langle\cdot,\cdot\rangle\) is the Euclidean dot product.

\section{Geometric reformulation of GC via \(\iota\)}

Because \(\iota\) is injective, the condition \(\eta(s)=0\) is equivalent to the spiral closing in the horizontal plane (the infinite sum of horizontal projections is zero). Using the spiral, we can rewrite the dot product:
\[
\langle \mathbf{S}_{N-1},\; \mathbf{v}_N \rangle = \left\langle -\sum_{n=N}^{\infty} \mathbf{v}_n,\; \mathbf{v}_N \right\rangle = -|\mathbf{v}_N|^2 - \left\langle \sum_{n=N+1}^{\infty} \mathbf{v}_n,\; \mathbf{v}_N \right\rangle.
\]
Thus GC becomes
\[
\left\langle \sum_{n=N+1}^{\infty} \mathbf{v}_n,\; \mathbf{v}_N \right\rangle \ge -|\mathbf{v}_N|^2.
\]
This is always true if the tail has a non‑negative projection onto \(\mathbf{v}_N\) up to the magnitude of \(\mathbf{v}_N\). But the tail itself is a sum of vectors that are nearly opposite to \(\mathbf{v}_N\) because \(\mathbf{v}_{N+1} \approx -\mathbf{v}_N\). Therefore the left side is approximately \(-|\mathbf{v}_N|^2\) plus smaller corrections. Whether it is \(\ge -|\mathbf{v}_N|^2\) depends on the exact phases and decay rates.

\section{Attempt to prove GC from \(\eta(s)=0\) using \(\iota\)}

Suppose \(\eta(s)=0\). Then the spiral satisfies \(\sum_{n=1}^{\infty} \mathbf{v}_n = \mathbf{0}\). For any \(N\), the vector \(\mathbf{S}_{N-1}\) is the negative of the tail \(T_N = \sum_{n=N}^{\infty} \mathbf{v}_n\). Hence
\[
\langle \mathbf{S}_{N-1},\mathbf{v}_N \rangle = -\langle T_N,\mathbf{v}_N \rangle = -|\mathbf{v}_N|^2 - \langle T_{N+1},\mathbf{v}_N \rangle.
\]
Thus GC is equivalent to
\[
\langle T_{N+1},\mathbf{v}_N \rangle \ge -|\mathbf{v}_N|^2. \tag{*}
\]
Now \(T_{N+1}\) is a sum of vectors that are almost opposite to each other and to \(\mathbf{v}_N\). Using the quasi‑geometric ratio
\[
\mathbf{v}_{n+1} = -\left(\frac{n}{n+1}\right)^{\sigma} e^{-i t\ln(1+1/n)} \mathbf{v}_n,
\]
we can write
\[
T_{N+1} = \mathbf{v}_{N+1} \sum_{k=0}^{\infty} \prod_{j=1}^{k} \rho_{N+j},\quad
\rho_{m} = -\left(\frac{m}{m+1}\right)^{\sigma} e^{-i t\ln(1+1/m)}.
\]
The infinite product converges because \(|\rho_m| \to 1\) and the arguments approach \(\pi\) from below. One can then compute \(\langle T_{N+1},\mathbf{v}_N \rangle = |\mathbf{v}_N|^2 \cdot \operatorname{Re}\bigl( \rho_N (1 + \rho_{N+1} + \rho_{N+1}\rho_{N+2} + \cdots ) \bigr)\). The condition (*) becomes
\[
\operatorname{Re}\bigl( \rho_N (1 + \rho_{N+1} + \rho_{N+1}\rho_{N+2} + \cdots ) \bigr) \ge -1.
\]
This is a purely numerical inequality that depends only on \(\sigma\) and \(t\) through \(\rho_m\). For a zero of \(\eta(s)\), the value of the infinite sum is determined by the condition that the total sum is zero. In fact, the requirement \(\sum_{n=1}^{\infty} \mathbf{v}_n = 0\) forces a specific relation between the \(\rho_m\) that might imply the inequality. However, deriving that relation explicitly is equivalent to solving for \(s\) such that the infinite alternating geometric‑like series sums to zero, which is exactly the original equation \(\eta(s)=0\). No simplification has been achieved.

\section{Why this does not prove GC unconditionally}

The above manipulation shows that GC is equivalent to a certain inequality involving the tail. But proving that inequality for all \(N\) given only that the total sum is zero is not easier than the original problem. In fact, the condition \(\eta(s)=0\) is a single equation in \(s\); it does not directly control each partial sum individually. The greedy condition is a set of infinitely many inequalities. For a zero not on the critical line, we have already shown (in the proof of Theorem 4.1, reverse direction) that GC fails asymptotically for large even \(N\). Therefore, any zero that satisfies GC must be on the line. Conversely, zeros on the line are known to satisfy GC via the Riemann–Siegel formula. Thus the equivalence is established, but the unconditional verification of GC for all zeros is exactly RH.

\section{Conclusion}

The iota function provides a geometric language to restate the greedy condition, but it does not offer a new proof. The statement “every zero satisfies GC” is equivalent to RH. The forward direction (RH ⇒ GC) is known, the reverse direction (GC ⇒ RH) is proven, but neither direction gives an independent verification of GC. Therefore, the Riemann Hypothesis remains open.

\begin{thebibliography}{9}
\bibitem{titchmarsh} E. C. Titchmarsh, \textit{The Theory of the Riemann Zeta Function}, Oxford University Press, 1986.
\bibitem{edwards} H. M. Edwards, \textit{Riemann's Zeta Function}, Dover, 2001.
\end{thebibliography}

\end{document}